\documentclass[11pt]{article}
\usepackage{amsmath, amssymb}
\usepackage{geometry}
\usepackage{array}
\usepackage{booktabs}
\geometry{margin=0.7in}
\setlength{\parindent}{0pt}
\setlength{\parskip}{0.35em}
\renewcommand{\arraystretch}{1.15}

\newcommand{\cheatsheettitle}[2]{%
\begin{center}
{\LARGE #1}\\[0.4em]
{\large #2}
\end{center}
}


\begin{document}
\cheatsheettitle{Algebra II Quadratics and Optimization Cheatsheet}{Module 3}

\section*{Forms}
Standard:
\[
f(x)=ax^2+bx+c
\]
Vertex:
\[
f(x)=a(x-h)^2+k
\]
Factored:
\[
f(x)=a(x-r_1)(x-r_2)
\]

\section*{Vertex}
From standard form:
\[
x=-\frac{b}{2a}
\]
\[
y=f\left(-\frac{b}{2a}\right)
\]
From vertex form:
\[
(h,k)
\]
Axis of symmetry:
\[
x=h
\]

\section*{Solving}
Factoring:
\[
AB=0\Rightarrow A=0\text{ or }B=0
\]
Completing the square:
\[
x^2+bx=\left(x+\frac b2\right)^2-\left(\frac b2\right)^2
\]
Quadratic formula:
\[
x=\frac{-b\pm\sqrt{b^2-4ac}}{2a}
\]

\section*{Discriminant}
\[
D=b^2-4ac
\]
\[
\begin{array}{ll}
D>0 & \text{two distinct real roots}\\
D=0 & \text{one repeated real root}\\
D<0 & \text{two complex conjugate roots}
\end{array}
\]

\section*{Complex Solutions}
\[
i^2=-1,\qquad \sqrt{-a}=i\sqrt a
\]
If coefficients are real and $a+bi$ is a root, then $a-bi$ is also a root.

\section*{Optimization}
For a quadratic model:
\[
f(x)=ax^2+bx+c
\]
The maximum or minimum occurs at:
\[
x=-\frac{b}{2a}
\]
\[
a>0\Rightarrow \text{minimum},\qquad a<0\Rightarrow \text{maximum}
\]
Revenue form:
\[
R(x)=x\cdot p(x)
\]

\section*{Quick Checks}
\begin{itemize}
  \item Choose the form that reveals the feature needed.
  \item Vertex form reveals vertex and transformations.
  \item Factored form reveals zeros.
  \item Standard form reveals \(y\)-intercept and supports the quadratic formula.
\end{itemize}

\end{document}
